% GENERATED FILE. Edit canonical lessons, then run npm run generate:books.
% canonical-source-hash: fnv1a64:07d1bff859438725
% canonical-lessons: UR-C438-chonch, UR-C438-pankh, UR-C438-ashiyana, UR-C438-chugna, UR-C438-shikar

\chapter{Beaks and Nests}
\label{ch:ur-beaks-and-nests}
\hlchaptermodality{438}

\begin{chapteropening}
\textbf{By the end of this chapter:} \emph{I can describe birds: beaks, feathers, nests and pecking.}

Complete the last lesson of chapter 438: I can describe birds: beaks, feathers, nests and pecking.
\end{chapteropening}

\section[\texorpdfstring{\ur{چونچ} (choṅch)}{چونچ (choṅch)}]{\ur{چونچ} (choṅch) — a beak}

\label{lesson:UR-C438-chonch}



\begin{quote}
Before the new one: say the Urdu for moss, then the Urdu for to come up (of plants).
\end{quote}

\subsection*{You'll want to know: \ur{چونچ}}
\textbf{\ur{چونچ}} — \emph{choṅch} — \textquotedblleft{}a beak\textquotedblright{}.

\begin{quote}
\ur{پرندے} \ur{کی} \ur{چونچ} \ur{میں} \ur{ایک} \ur{تنکا} \ur{ہے۔}
\end{quote}

\begin{quote}
\emph{parinde kī choṅch meṅ ek tinkā hai.}
\end{quote}

\begin{quote}
The bird has a straw in its beak.
\end{quote}

Feminine.

\begin{cousinweb}
\textbf{\ur{چونچ}} (\emph{choṅch}) is inherited: CDIAL lists the Hindi cousin \emph{cõc}, \textquotedblleft{}a beak\textquotedblright{}, under Sanskrit \emph{cañcu} (CDIAL 4569).
\end{cousinweb}

\begin{grammarlens}[title={What you've built}]
One more word for getting around.
\end{grammarlens}

\subsection*{Your turn}
\begin{itemize}
\raggedright
  \item \emph{Say it:} \emph{choṅch}
  \item \emph{Say it:} \emph{choṅch}, once more
  \item \emph{Say it:} \emph{parinde kī choṅch meṅ ek tinkā hai.} — The bird has a straw in its beak.
  \item \emph{Recall:} say the Urdu for moss, then the Urdu for to come up (of plants), then say \emph{choṅch} again
\end{itemize}

\begin{tcolorbox}[breakable,colback=teal!4,colframe=teal!35!black,title={Before you move on}]
What does \ur{چونچ} mean? (\textquotedblleft{}A beak\textquotedblright{}.) Say it once more.
\end{tcolorbox}
\section[\texorpdfstring{\ur{پنکھ} (pankh)}{پنکھ (pankh)}]{\ur{پنکھ} (pankh) — a wing; a feather}

\label{lesson:UR-C438-pankh}



\begin{quote}
Before the new one: say the Urdu for to come up (of plants), then the Urdu for a beak.
\end{quote}

\subsection*{You'll want to know: \ur{پنکھ}}
\textbf{\ur{پنکھ}} — \emph{pankh} — \textquotedblleft{}a wing; a feather\textquotedblright{}.

\begin{quote}
\ur{طوطے} \ur{کے} \ur{پنکھ} \ur{ہرے} \ur{ہیں۔}
\end{quote}

\begin{quote}
\emph{tote ke pankh hare haiṅ.}
\end{quote}

\begin{quote}
The parrot's feathers are green.
\end{quote}

Masculine.

\begin{cousinweb}
\textbf{\ur{پنکھ}} (\emph{pankh}) is inherited from Sanskrit \emph{pakṣá}, \textquotedblleft{}a wing, a feather\textquotedblright{}, under which CDIAL lists the Hindi cousin \emph{pā̃kh}, \textquotedblleft{}a wing\textquotedblright{} (CDIAL 7627); \textbf{\ur{پنکھا}}, \textquotedblleft{}a fan\textquotedblright{}, comes from the same word.
\end{cousinweb}

\begin{grammarlens}[title={What you've built}]
One more word for getting around.
\end{grammarlens}

\subsection*{Your turn}
\begin{itemize}
\raggedright
  \item \emph{Say it:} \emph{pankh}
  \item \emph{Say it:} \emph{pankh}, once more
  \item \emph{Say it:} \emph{tote ke pankh hare haiṅ.} — The parrot's feathers are green.
  \item \emph{Recall:} say the Urdu for to come up (of plants), then the Urdu for a beak, then say \emph{pankh} again
\end{itemize}

\begin{tcolorbox}[breakable,colback=teal!4,colframe=teal!35!black,title={Before you move on}]
What does \ur{پنکھ} mean? (\textquotedblleft{}A wing; a feather\textquotedblright{}.) Say it once more.
\end{tcolorbox}
\section[\texorpdfstring{\ur{آشیانہ} (āshiyānā)}{آشیانہ (āshiyānā)}]{\ur{آشیانہ} (āshiyānā) — a nest}

\label{lesson:UR-C438-ashiyana}



\begin{quote}
Before the new one: say the Urdu for a beak, then the Urdu for a wing.
\end{quote}

\subsection*{You'll want to know: \ur{آشیانہ}}
\textbf{\ur{آشیانہ}} — \emph{āshiyānā} — \textquotedblleft{}a nest\textquotedblright{}.

\begin{quote}
\ur{چڑیا} \ur{نے} \ur{درخت} \ur{پر} \ur{آشیانہ} \ur{بنایا۔}
\end{quote}

\begin{quote}
\emph{chiṛyā ne darakht par āshiyānā banāyā.}
\end{quote}

\begin{quote}
The sparrow built a nest in the tree.
\end{quote}

Masculine.

\begin{cousinweb}
\textbf{\ur{آشیانہ}} (\emph{āshiyānā}) is usually traced to Persian \emph{āshiyāna}, \textquotedblleft{}a nest\textquotedblright{}; poets also use it for a home.
\end{cousinweb}

\begin{grammarlens}[title={What you've built}]
One more word for getting around.
\end{grammarlens}

\subsection*{Your turn}
\begin{itemize}
\raggedright
  \item \emph{Say it:} \emph{āshiyānā}
  \item \emph{Say it:} \emph{āshiyānā}, once more
  \item \emph{Say it:} \emph{chiṛyā ne darakht par āshiyānā banāyā.} — The sparrow built a nest in the tree.
  \item \emph{Recall:} say the Urdu for a beak, then the Urdu for a wing, then say \emph{āshiyānā} again
\end{itemize}

\begin{tcolorbox}[breakable,colback=teal!4,colframe=teal!35!black,title={Before you move on}]
What does \ur{آشیانہ} mean? (\textquotedblleft{}A nest\textquotedblright{}.) Say it once more.
\end{tcolorbox}
\section[\texorpdfstring{\ur{چگنا} (chugnā)}{چگنا (chugnā)}]{\ur{چگنا} (chugnā) — to peck (at); to pick up (grain)}

\label{lesson:UR-C438-chugna}



\begin{quote}
Before the new one: say the Urdu for a wing, then the Urdu for a nest.
\end{quote}

\subsection*{You'll want to know: \ur{چگنا}}
\textbf{\ur{چگنا}} — \emph{chugnā} — \textquotedblleft{}to peck (at); to pick up (grain)\textquotedblright{}.

\begin{quote}
\ur{کبوتر} \ur{زمین} \ur{سے} \ur{چاول} \ur{چگتے} \ur{ہیں۔}
\end{quote}

\begin{quote}
\emph{kabūtar zamīn se chāval chugte haiṅ.}
\end{quote}

\begin{quote}
The pigeons peck rice off the ground.
\end{quote}

Past, with \ur{نے} (he or she …): \ur{اس} \ur{نے} \ur{چاول} \ur{چگے}, \emph{us ne chāval chuge}.

\begin{cousinweb}
\textbf{\ur{چگنا}} (\emph{chugnā}) is inherited: CDIAL lists the Hindi cousin \emph{cugnā}, \textquotedblleft{}to peck, to pick up\textquotedblright{}, under \emph{cugyati}, a form it reconstructs (CDIAL 4852).
\end{cousinweb}

\begin{grammarlens}[title={What you've built}]
One more word for getting around.
\end{grammarlens}

\subsection*{Your turn}
\begin{itemize}
\raggedright
  \item \emph{Say it:} \emph{chugnā}
  \item \emph{Say it:} \emph{chugnā}, once more
  \item \emph{Say it:} \emph{kabūtar zamīn se chāval chugte haiṅ.} — The pigeons peck rice off the ground.
  \item \emph{Recall:} say the Urdu for a wing, then the Urdu for a nest, then say \emph{chugnā} again
\end{itemize}

\begin{tcolorbox}[breakable,colback=teal!4,colframe=teal!35!black,title={Before you move on}]
What does \ur{چگنا} mean? (\textquotedblleft{}To peck (at); to pick up (grain)\textquotedblright{}.) Say it once more.
\end{tcolorbox}
\section[\texorpdfstring{\ur{شکار} (shikār)}{شکار (shikār)}]{\ur{شکار} (shikār) — prey; hunting}

\label{lesson:UR-C438-shikar}



\begin{quote}
Before the new one: say the Urdu for a nest, then the Urdu for to peck (at).
\end{quote}

\subsection*{You'll want to know: \ur{شکار}}
\textbf{\ur{شکار}} — \emph{shikār} — \textquotedblleft{}prey; hunting\textquotedblright{}.

\begin{quote}
\ur{بلی} \ur{کا} \ur{شکار} \ur{ایک} \ur{چوہا} \ur{تھا۔}
\end{quote}

\begin{quote}
\emph{billī kā shikār ek chūhā thā.}
\end{quote}

\begin{quote}
The cat's prey was a mouse.
\end{quote}

Masculine.

\begin{cousinweb}
\textbf{\ur{شکار}} (\emph{shikār}) is usually traced to Persian \emph{shikār}, \textquotedblleft{}hunting, the chase\textquotedblright{}; in Urdu it is both the hunt and the animal hunted.
\end{cousinweb}

\begin{grammarlens}[title={What you've built}]
One more word for getting around.
\end{grammarlens}

\subsection*{Your turn}
\begin{itemize}
\raggedright
  \item \emph{Say it:} \emph{shikār}
  \item \emph{Say it:} \emph{shikār}, once more
  \item \emph{Say it:} \emph{billī kā shikār ek chūhā thā.} — The cat's prey was a mouse.
  \item \emph{Recall:} say the Urdu for a nest, then the Urdu for to peck (at), then say \emph{shikār} again
\end{itemize}

\begin{tcolorbox}[breakable,colback=teal!4,colframe=teal!35!black,title={Before you move on}]
What does \ur{شکار} mean? (\textquotedblleft{}Prey; hunting\textquotedblright{}.) Say it once more.
\end{tcolorbox}
